\label{SI:priormean}
{\color{blue}
To assess the sensitivity of the EnKF to the choice of prior mean parameters, all nominal parameter values from Kotidis \textit{et al.} \cite{Kotidis2019Model-basedCulture} were simultaneously scaled by $\pm10\%$ and $\pm20\%$, and the filter was re-run for all six datasets. Figures~\ref{fig:priormean_A_flask}--\ref{fig:priormean_B_pl40} compare the resulting state trajectories to the nominal baseline and experimental data. Across most datasets and states, the EnKF trajectories remain qualitatively consistent across the perturbation range, indicating that the filter is broadly robust to moderate uncertainty in the prior mean parameter values. Minor differences in transient dynamics are observed for certain metabolite states, most notably in the Cell Line B datasets where the system deviates most substantially from the calibration conditions. The sole exception is asparagine (Asn) in the Cell Line B Feed C and Feed U datasets at the $-20\%$ perturbation level, which exceeds the 50\% deviation threshold. Table~\ref{tab:priormean_stability} summarises per-state stability across all perturbation levels using the same criterion as Table~\ref{tab:priorcov_stability}.
}

\begin{figure}[H]
    \centering
    \includegraphics[width=0.95\textwidth]{Figs/param_sensitivity_CHO_T127_flask_PMJ.png}
    \caption{\rev{Effect of $\pm10\%$ and $\pm20\%$ perturbations to all prior mean parameters on EnKF state estimation for Cell Line A Shake Flask. Trajectories are shown for $-20\%$ (dashed, darkest), $-10\%$ (solid, dark), $+10\%$ (solid, light), and $+20\%$ (dotted, lightest) perturbations, alongside experimental measurements (markers, $\pm$standard deviation). The baseline ($0\%$) corresponds to the nominal parameter values from Kotidis \textit{et al.}~\cite{Kotidis2019Model-basedCulture}.}}
    \label{fig:priormean_A_flask}
\end{figure}

\begin{figure}[H]
    \centering
    \includegraphics[width=0.95\textwidth]{Figs/param_sensitivity_CHO_T127_SNS_36.5.png}
    \caption{\rev{Effect of $\pm10\%$ and $\pm20\%$ perturbations to all prior mean parameters on EnKF state estimation for Cell Line A Bioreactor 36.5\textdegree C. Trajectories are shown for $-20\%$ (dashed, darkest), $-10\%$ (solid, dark), $+10\%$ (solid, light), and $+20\%$ (dotted, lightest) perturbations, alongside experimental measurements (markers, $\pm$standard deviation). The baseline ($0\%$) corresponds to the nominal parameter values from Kotidis \textit{et al.}~\cite{Kotidis2019Model-basedCulture}.}}
    \label{fig:priormean_A_36}
\end{figure}

\begin{figure}[H]
    \centering
    \includegraphics[width=0.95\textwidth]{Figs/param_sensitivity_CHO_T127_SNS_32.png}
    \caption{\rev{Effect of $\pm10\%$ and $\pm20\%$ perturbations to all prior mean parameters on EnKF state estimation for Cell Line A Bioreactor 32\textdegree C. Trajectories are shown for $-20\%$ (dashed, darkest), $-10\%$ (solid, dark), $+10\%$ (solid, light), and $+20\%$ (dotted, lightest) perturbations, alongside experimental measurements (markers, $\pm$standard deviation). The baseline ($0\%$) corresponds to the nominal parameter values from Kotidis \textit{et al.}~\cite{Kotidis2019Model-basedCulture}.}}
    \label{fig:priormean_A_32}
\end{figure}

\begin{figure}[H]
    \centering
    \includegraphics[width=0.95\textwidth]{Figs/param_sensitivity_CHO_GS46_F_C_Inv.png}
    \caption{\rev{Effect of $\pm10\%$ and $\pm20\%$ perturbations to all prior mean parameters on EnKF state estimation for Cell Line B Feed C. Trajectories are shown for $-20\%$ (dashed, darkest), $-10\%$ (solid, dark), $+10\%$ (solid, light), and $+20\%$ (dotted, lightest) perturbations, alongside experimental measurements (markers, $\pm$standard deviation). The baseline ($0\%$) corresponds to the nominal parameter values from Kotidis \textit{et al.}~\cite{Kotidis2019Model-basedCulture}.}}
    \label{fig:priormean_B_C}
\end{figure}

\begin{figure}[H]
    \centering
    \includegraphics[width=0.95\textwidth]{Figs/param_sensitivity_CHO_GS46_F_all.png}
    \caption{\rev{Effect of $\pm10\%$ and $\pm20\%$ perturbations to all prior mean parameters on EnKF state estimation for Cell Line B Feed U. Trajectories are shown for $-20\%$ (dashed, darkest), $-10\%$ (solid, dark), $+10\%$ (solid, light), and $+20\%$ (dotted, lightest) perturbations, alongside experimental measurements (markers, $\pm$standard deviation). The baseline ($0\%$) corresponds to the nominal parameter values from Kotidis \textit{et al.}~\cite{Kotidis2019Model-basedCulture}.}}
    \label{fig:priormean_B_U}
\end{figure}

\begin{figure}[H]
    \centering
    \includegraphics[width=0.95\textwidth]{Figs/param_sensitivity_CHO_GS46_F_all_pl40.png}
    \caption{\rev{Effect of $\pm10\%$ and $\pm20\%$ perturbations to all prior mean parameters on EnKF state estimation for Cell Line B Feed U plus 40\%. Trajectories are shown for $-20\%$ (dashed, darkest), $-10\%$ (solid, dark), $+10\%$ (solid, light), and $+20\%$ (dotted, lightest) perturbations, alongside experimental measurements (markers, $\pm$standard deviation). The baseline ($0\%$) corresponds to the nominal parameter values from Kotidis \textit{et al.}~\cite{Kotidis2019Model-basedCulture}.}}
    \label{fig:priormean_B_pl40}
\end{figure}

% cell-colour shorthands (local to this file)
\definecolor{stabgreen}{RGB}{163,210,163}   % muted sage green — stable
\definecolor{unstabred}{RGB}{218,112,112}   % muted crimson — unstable
\providecommand{\G}{\cellcolor{stabgreen}\ding{51}}   % tick
\providecommand{\R}{\cellcolor{unstabred}\ding{55}}   % cross

\begin{table}[ht]
\centering
\renewcommand{\arraystretch}{1.3}
\small
\setlength{\tabcolsep}{3pt}
\caption{\rev{Per-state EnKF stability under different prior mean parameter perturbations for all six datasets. Green cells indicate a stable trajectory; red cells indicate that the trajectory deviates from the $0\%$ baseline by more than 50\% of the experimental maximum for that state at any simulation point. The $0\%$ column is the nominal case used throughout the paper and is stable by definition. Asparagine in Cell Line B Feed C and Feed U is flagged unstable at $-20\%$.}}
\label{tab:priormean_stability}
\arrayrulecolor{black}
\adjustbox{max width=\textwidth}{%
\begin{tabular}{|l|p{1.6cm}|>{\centering\arraybackslash}p{2.2cm}|>{\centering\arraybackslash}p{2.2cm}|>{\centering\arraybackslash}p{2.2cm}|>{\centering\arraybackslash}p{2.2cm}|>{\centering\arraybackslash}p{2.2cm}|}
\hline
\rowcolor{gray!30}
\textbf{Dataset} & \textbf{State} &
\textbf{\makecell{$-20\%$\\Prior Mean}} &
\textbf{\makecell{$-10\%$\\Prior Mean}} &
\textbf{\makecell{$0\%$\\Prior Mean\\(Baseline)}} &
\textbf{\makecell{$+10\%$\\Prior Mean}} &
\textbf{\makecell{$+20\%$\\Prior Mean}} \\
\hline
%──────────────────────────────────────────────────────
% Cell Line A — Shake Flask
%──────────────────────────────────────────────────────
\multirow{8}{*}{\makecell[l]{Cell Line A\\Shake Flask}}
 & Xv  & \G & \G & \G & \G & \G \\ \hhline{|~|------|}
 & mAb & \G & \G & \G & \G & \G \\ \hhline{|~|------|}
 & Glc & \G & \G & \G & \G & \G \\ \hhline{|~|------|}
 & Amm & \G & \G & \G & \G & \G \\ \hhline{|~|------|}
 & Gln & \G & \G & \G & \G & \G \\ \hhline{|~|------|}
 & Lac & \G & \G & \G & \G & \G \\ \hhline{|~|------|}
 & Glu & \G & \G & \G & \G & \G \\ \hhline{|~|------|}
 & Asn & \G & \G & \G & \G & \G \\
\hline
%──────────────────────────────────────────────────────
% Cell Line A — Bioreactor 36.5 °C
%──────────────────────────────────────────────────────
\multirow{8}{*}{\makecell[l]{Cell Line A\\Bioreactor 36.5\textdegree C}}
 & Xv  & \G & \G & \G & \G & \G \\ \hhline{|~|------|}
 & mAb & \G & \G & \G & \G & \G \\ \hhline{|~|------|}
 & Glc & \G & \G & \G & \G & \G \\ \hhline{|~|------|}
 & Amm & \G & \G & \G & \G & \G \\ \hhline{|~|------|}
 & Gln & \G & \G & \G & \G & \G \\ \hhline{|~|------|}
 & Lac & \G & \G & \G & \G & \G \\ \hhline{|~|------|}
 & Glu & \G & \G & \G & \G & \G \\ \hhline{|~|------|}
 & Asn & \G & \G & \G & \G & \G \\
\hline
%──────────────────────────────────────────────────────
% Cell Line A — Bioreactor 32 °C
%──────────────────────────────────────────────────────
\multirow{8}{*}{\makecell[l]{Cell Line A\\Bioreactor 32\textdegree C}}
 & Xv  & \G & \G & \G & \G & \G \\ \hhline{|~|------|}
 & mAb & \G & \G & \G & \G & \G \\ \hhline{|~|------|}
 & Glc & \G & \G & \G & \G & \G \\ \hhline{|~|------|}
 & Amm & \G & \G & \G & \G & \G \\ \hhline{|~|------|}
 & Gln & \G & \G & \G & \G & \G \\ \hhline{|~|------|}
 & Lac & \G & \G & \G & \G & \G \\ \hhline{|~|------|}
 & Glu & \G & \G & \G & \G & \G \\ \hhline{|~|------|}
 & Asn & \G & \G & \G & \G & \G \\
\hline
%──────────────────────────────────────────────────────
% Cell Line B — Feed C
%──────────────────────────────────────────────────────
\multirow{8}{*}{\makecell[l]{Cell Line B\\Feed C}}
 & Xv  & \G & \G & \G & \G & \G \\ \hhline{|~|------|}
 & mAb & \G & \G & \G & \G & \G \\ \hhline{|~|------|}
 & Glc & \G & \G & \G & \G & \G \\ \hhline{|~|------|}
 & Amm & \G & \G & \G & \G & \G \\ \hhline{|~|------|}
 & Gln & \G & \G & \G & \G & \G \\ \hhline{|~|------|}
 & Lac & \G & \G & \G & \G & \G \\ \hhline{|~|------|}
 & Glu & \G & \G & \G & \G & \G \\ \hhline{|~|------|}
 & Asn & \R & \G & \G & \G & \G \\
\hline
%──────────────────────────────────────────────────────
% Cell Line B — Feed U
%──────────────────────────────────────────────────────
\multirow{8}{*}{\makecell[l]{Cell Line B\\Feed U}}
 & Xv  & \G & \G & \G & \G & \G \\ \hhline{|~|------|}
 & mAb & \G & \G & \G & \G & \G \\ \hhline{|~|------|}
 & Glc & \G & \G & \G & \G & \G \\ \hhline{|~|------|}
 & Amm & \G & \G & \G & \G & \G \\ \hhline{|~|------|}
 & Gln & \G & \G & \G & \G & \G \\ \hhline{|~|------|}
 & Lac & \G & \G & \G & \G & \G \\ \hhline{|~|------|}
 & Glu & \G & \G & \G & \G & \G \\ \hhline{|~|------|}
 & Asn & \R & \G & \G & \G & \G \\
\hline
%──────────────────────────────────────────────────────
% Cell Line B — Feed U plus 40%
%──────────────────────────────────────────────────────
\multirow{8}{*}{\makecell[l]{Cell Line B\\Feed U plus 40\%}}
 & Xv  & \G & \G & \G & \G & \G \\ \hhline{|~|------|}
 & mAb & \G & \G & \G & \G & \G \\ \hhline{|~|------|}
 & Glc & \G & \G & \G & \G & \G \\ \hhline{|~|------|}
 & Amm & \G & \G & \G & \G & \G \\ \hhline{|~|------|}
 & Gln & \G & \G & \G & \G & \G \\ \hhline{|~|------|}
 & Lac & \G & \G & \G & \G & \G \\ \hhline{|~|------|}
 & Glu & \G & \G & \G & \G & \G \\ \hhline{|~|------|}
 & Asn & \G & \G & \G & \G & \G \\
\hline
\end{tabular}}
\end{table}
